\documentclass{article} % For LaTeX2e
\usepackage{iclr2027_conference,times}

\input{math_commands.tex}

\usepackage{hyperref}
\usepackage{url}
\usepackage{graphicx}
\usepackage{booktabs}
\usepackage{amsmath,amssymb}

\title{Tree-Like Class Geometry in Foundation Models: \\ A Zero-Cost Diagnostic for Metric Selection}

% Anonymous while \iclrfinalcopy is commented out.
\author{Anonymous authors}

%\iclrfinalcopy % Uncomment for camera-ready version, but NOT for submission.

\begin{document}

\maketitle

\begin{abstract}
% SKELETON (v0). Three validated statements, in this order:
% (1) FORM: tree-like inter-class metric structure is near-universal across
%     foundation models (69/72 model-dataset cells beyond nulls matched in
%     dimension and covariance spectrum), and it is the training objective,
%     not the modality, that determines its strength.
% (2) CONTENT: the tree each model builds is its own; tree-likeness and
%     WordNet alignment decouple (label-free DINOv2 is the most tree-like
%     and the least aligned), and hyperbolic distances do not increase
%     cross-model agreement. Convergence in form, not in content.
% (3) TOOL: two zero-cost descriptors (delta-hyperbolicity vs matched null,
%     ORC) decide WHEN to leave Euclidean distances, and the training
%     objective selects WHICH metric collects the gain (angular for SSL,
%     Poincare for contrastive VLMs); two-step protocol, no training.
The geometry of learned representations bears directly on when non-Euclidean distance functions help downstream. We study the metric structure of class representations, the inter-class distance matrices of frozen models, across 30 foundation models: 14 vision backbones, 9 causal language models, and 7 text embedders. Tree-likeness is measured with two zero-cost, coordinate-free diagnostics, Gromov $\delta$-hyperbolicity and Ollivier--Ricci curvature, calibrated on reference geometries and, crucially, read against nulls matched in dimension and covariance spectrum, since random high-dimensional clouds concentrate toward degenerate star trees and make raw $\delta$ uninterpretable on its own. Three findings emerge. First, the \emph{form} is near-universal: 69 of 72 vision model$\times$dataset combinations show tree structure beyond their matched null; its strength is determined by the training objective rather than the modality, deepens with model scale, and a parallelogram curvature-sign estimator confirms increasingly negative curvature with scale for self-supervised models. Second, the \emph{content} is model-specific: the most tree-like family (label-free DINOv2) is the least aligned with WordNet, and hyperbolic distances do not increase cross-model agreement; models converge in the form of their class geometry, not in its content. Third, the diagnostics support a training-free, two-step metric-selection protocol: $\delta$ and ORC predict when leaving Euclidean distances pays on prototype-based tasks, and the training objective selects which zero-cost metric collects the gain, cosine for self-supervised models, a parameter-free Poincar\'{e} projection for contrastive vision--language models ($+0.9$ to $+1.3$pp over cosine after L2 normalization, up to $+2.05$pp over Euclidean, McNemar $p \approx 4\times10^{-38}$). On multi-level caption hierarchies the radial axis of the projection aligns with ground-truth semantic generality, and the protocol transfers to text, correctly predicting when hyperbolic distances help little. We also report where the projection does not help, sample-level kNN and retrieval, scoping its domain of validity to prototype-based metric tasks.
\end{abstract}

\section{Introduction}
% NEW FRAMING: within-model claim from the title down. PRH appears only as a
% one-sentence pointer to related work. Every number below traces to
% table1_regenerated.csv / exp CSVs (see comments).

Semantic knowledge is hierarchical: classes nest into superclasses, concepts into broader concepts. When a frozen foundation model represents a set of classes, as image-class centroids or as embedded class descriptions, the pairwise distances between those representations form a metric space, and the structure of that space decides which distance functions serve downstream tasks well. If the inter-class metric structure is tree-like, distances respecting negative curvature or angular geometry can outperform raw Euclidean distances at zero training cost. Most work in hyperbolic deep learning \emph{imposes} hierarchical structure through specialized architectures, losses, or trainable Poincar\'{e} embeddings \citep{nickel2017poincare, Khrulkov_2020_CVPR, chen2022fully, desai2023meru}. We ask the converse, purely observational question: \emph{do frozen foundation models already organize class concepts in a tree-like metric structure, and when is that structure exploitable?}

Answering this requires care, because the na\"ive measurement is confounded. Gromov $\delta$-hyperbolicity \citep{Gromov1987} assigns $0$ to trees, but random high-dimensional point clouds concentrate toward equidistance, a degenerate star tree, so a low normalized $\delta$ alone is not evidence of learned hierarchy; and the reference constants themselves depend on which definition of hyperbolicity is used (the four-point constant of the hyperbolic plane is $\ln 2$, not the thin-triangles constant $\ln(1+\sqrt{2})$). % \citep{nica2016strong} TODO-bib
Our measurement layer therefore rests on three legs: calibration of the exact estimator on reference geometries (trees, hyperbolic planes, spheres, matched Gaussians), \emph{excess over a null matched in dimension and covariance spectrum} as the primary reported quantity, and a direct parallelogram estimator of the curvature \emph{sign}. Together with Ollivier--Ricci curvature \citep{ollivier2009ricci} as a complementary local descriptor, this turns ``how tree-like is this representation?'' into a calibrated, falsifiable question.

Three findings emerge, which the title condenses. \textbf{Form is near-universal:} 69 of 72 vision model$\times$dataset combinations exhibit tree structure beyond their matched null, the exceptions being self-supervised DINOv2 on its pretraining-like distribution. % exp11_null_per_dataset.csv
The strength of the structure is set by the training objective rather than the modality, deepens with model scale on transfer data (DINOv2-S$\to$G excess $-0.087\to-0.146$ on CIFAR-10), and the curvature sign turns increasingly negative with scale for the same family. % exp12_curvature_sign.csv
Ablations tie the structure to learning: it vanishes under random weights, weakens under non-hierarchical fine-tuning, and emerges progressively across layers. \textbf{Content is model-specific:} the most tree-like family (label-free DINOv2) is the \emph{least} aligned with the WordNet taxonomy, ground-truth checks (WordNet correlation, superclass recovery, definition-only prompts, multi-level caption hierarchies) separate ``how tree-like'' from ``whose tree'', and hyperbolic distances do not increase cross-model agreement. Models converge in the form of their class geometry, not in its content; we discuss the relation to representational-convergence hypotheses in Section~2. \textbf{The geometry is a usable signal:} the two diagnostics, computable in seconds on training centroids, predict when non-Euclidean metrics beat Euclidean distances on prototype-based tasks, and the training objective selects \emph{which} zero-cost metric collects the gain: cosine suffices for self-supervised models, while for contrastive vision--language models a parameter-free Poincar\'{e} projection beats cosine even after L2 normalization. % exp2_metric_controls.csv, exp2b_normalized_stack.csv
We validate the resulting two-step protocol with paired statistics, transfer it to a text hierarchy where it correctly predicts marginal gains, and report the negative results that scope it: sample-level kNN and retrieval do not benefit, because the hierarchy lives at the prototype level while the radial map distorts local neighborhoods.

Our contributions are:
\begin{itemize}
\item \textbf{A calibrated measurement methodology} for tree-likeness of inter-class metric structure: estimator calibration on reference geometries, spectrum-matched nulls as the zero point, and a curvature-sign estimator, addressing the dimension confound that affects raw $\delta$-hyperbolicity readings.
\item \textbf{A geometric census of 30 foundation models} establishing that tree-like form is near-universal (69/72 cells beyond nulls), objective-determined, scale-deepening, and learned, while its content is model-specific: convergence in form, not in content.
\item \textbf{Semantic-correctness validation} against ground truth: WordNet alignment ($\rho$ up to $+0.59$), CIFAR-100 superclass recovery (ARI up to $0.61$), robustness to prompt templates including definition-only prompts with no class names, and radial alignment with semantic generality on multi-level caption hierarchies (strictly monotone chains up to $54\%$ vs.\ $4.2\%$ chance). % exp3, exp7, exp17, exp16
\item \textbf{A training-free metric-selection protocol}: when to leave Euclidean distances (predicted by $\delta$/ORC) and which zero-cost metric to use (set by the training objective), with paired confidence intervals and McNemar tests, a text-modality transfer, and explicitly scoped negative results.
\end{itemize}

\section{Related Work and Background}
% COMMITTED (1Eqj): 2.3 becomes Background (preliminaries); PRH paragraph
% merged into related work with the critical literature (Koepke et al. 2026;
% Groger et al. 2026 "Aristotelian View") and our form-not-content position
% promoted from old §2.4.
TODO

\section{Measuring Tree-Likeness}
% COMMITTED (1Eqj/RJje):
% - motivate both descriptors (delta: global, coordinate-free, embeddability;
%   ORC: local, edge-level, complementary).
% - CALIBRATION TABLE (from rebuttal): tree 0.000; H2 regions .162/.087/.043/.022;
%   S99 .143 chord / .179 geodesic; iid Gaussian .104/.061/.046.
% - CONSTANTS FIXED: four-point delta of H2 is ln 2 ~ 0.693 (Nica & Spakula 2016),
%   thin-triangles constant ln(1+sqrt2) ~ 0.881 is a DIFFERENT definition;
%   absolute vs normalized scales never mixed in one sentence.
% - MATCHED NULLS ARE PRIMARY: excess over spectrum-matched null is THE quantity;
%   "low delta" defined as "below the matched null", never absolute.
% - Exponential-map subsection rewritten with geometric reading + diagram.
% - Few-shot protocol as explicit algorithm box.
TODO

\section{The Geometry: Form Is Near-Universal, Content Is Model-Specific}
% MAIN-TEXT SUMMARY TABLE (committed to 1Eqj, shown in rebuttal).
% - 69/72 cells beyond matched nulls; exceptions DINOv2-S/B/G on ImageNet;
%   capacity trend on the EXCESS basis (S->G CIFAR-10: -0.087 -> -0.146).
% - Sphere disfavored: sphericizing raises delta-hat (12/12 geodesic, 10/12 chord).
% - ORC bridge analysis (committed to RJje): within-cluster vs bridge edges.
% - WordNet alignment section (committed to pux6): rho up to +0.59, ARI 0.61,
%   the DINOv2 decoupling, no-pooling control, random-groupings control.
% - Prompt robustness (committed to Xbn5): embedders robust, causal-LM claims
%   narrowed to mean+-range across templates.
TODO

\section{The Tool: When to Leave Euclidean, and Which Metric Collects the Gain}
% TWO-STEP PROTOCOL (from the discussion, the pux6/Xbn5 replies):
% Step 1: delta-hat/ORC decide whether ANY non-Euclidean metric helps
%   (best-metric advantage correlates with delta-hat; on MNIST nothing helps).
% Step 2: objective selects the metric: SSL -> cosine; contrastive VLM ->
%   Poincare (+0.9..+1.3pp over cosine on transfer datasets after L2);
%   supervised -> either. RT control: the gain is the metric, not the map.
% - Statistics: per-episode paired CIs; McNemar (242 vs 37, p ~ 4e-38).
% - DBpedia text transfer: diagnostic correctly predicts no meaningful gain.
% - HONEST NEGATIVES in the main text: kNN -2.6pp, retrieval 0/60;
%   domain of validity = prototype-based metric tasks, with mechanism.
TODO

\section{Ablations: The Geometry Is Learned}
% Port from neurips_2026.tex (random weights, fine-tuning, layer-wise),
% restated as excesses where cross-dimension comparisons occur.
TODO

\section{Discussion and Limitations}
% - No curvature estimation (descriptive claims are metric tree-likeness).
% - Template dependence for causal LMs.
% - Class inventory is human-chosen; labels probe, not create, the structure.
% - HierarCaps-style multi-level evaluation (committed to pux6) if ready;
%   otherwise stated as ongoing.
TODO

\subsubsection*{Reproducibility Statement}
% Code + frozen features + per-cell CSVs; every number traceable to a script.
TODO

\subsubsection*{Ethics Statement}
TODO

\subsubsection*{The Use of Large Language Models}
% MANDATORY for ICLR 2027. Draft honestly per the AI policy.
TODO

\bibliography{iclr2027_conference}
\bibliographystyle{iclr2027_conference}

\appendix
\section{Full Per-Cell Tables}
% Table 1 REGENERATED from tasks_t707 (internal audit: NC-H column + the
% ImageNet delta cells DINOv2-G .066->.080, CLIP-L .108->.118). Split into
% per-section tables. Null tables, metric-control tables, prompt tables,
% McNemar per-cell, DBpedia.
TODO

\end{document}
